\documentclass[12pt, letterpaper]{article}
\usepackage{graphicx}

\newcommand{\ikjpg}[1]{
% \newpage
    \begin{center}
    % \centering
    \vspace{10pt}
    \includegraphics[width=\textwidth]{./jpg/Neandertal.0#1.jpeg}
    \vspace{10pt}
    %   \caption{Caption text.}
\end{center}
}

\begin{document}
\ikjpg{01}
Buenos días, Una de las preguntas que todos tenemos sobre la humanidad es: ¿qué es lo que hace que los humanos sean tan especiales?
¿Especiales en qué sentido? En el sentido del éxito que tenemos en este planeta. Entre los mamíferos somos miles de millones, más si incluimos nuestros animales, el ganado y los otros animales domésticos, hemos cubierto todo el planeta, todos los ecosistemas. Estamos en el Polo Norte, en el desierto australiano, y eso no es reciente hace mucho tiempo que los humanos han tenido un éxito inmenso.

Para ser más preciso, la pregunta que queremos explorar aquí es: ¿cuál es nuestra diferencia con los animales que hace que tengamos ese éxito?

\ikjpg{02}


Lo primero que nosotros pensamos es la inteligencia. Pensamos, bueno, sí, la diferencia, los humanos, lo que pasa es que los humanos son inteligentes, son mucho más inteligentes.

Pero en realidad hay que darse cuenta que no somos tan inteligentes, es decir, colectivamente somos muy inteligentes. Tenemos a Einstein, pero no todo el mundo es Einstein, pero Einstein no sabía de poesía. Y en realidad de manera colectiva sí tenemos toda esa inteligencia, pero de manera individual no somos mucho más inteligentes que por ejemplo estos mamíferos marinos, belugas, que están ahí haciendo anillos de burbuja. De manera colectiva, uno se da cuenta ahí que hay una inteligencia bastante grande y en realidad lo tienen los chimpancés, hay muchos animales que son inteligentes. Y esos animales no tienen mucho éxito, más bien estos están en peligro de extinción.

Entonces eso no explica nuestro éxito. La inteligencia no pareciera que es suficiente para explicar nuestro éxito. Sí tenemos más inteligencia, pero no mucho más que los demás animales y nada para explicar este éxito.
\ikjpg{03}
Entonces, una cosa que pensamos, otra idea que puede decir que nos diferencia de los animales y que eso sí parece explicar el éxito que podemos tener, es que somos seres de cultura. Para nosotros tenemos los genes, como los animales, pero además tenemos mucha información que viene como cultura.

Por ejemplo, si los esquimales pueden sobrevivir en el polo norte es porque tienen cultura. Por ejemplo, alguien que es un explorador europeo que fuera en el siglo XIX para el norte allá, él no podía sobrevivir y muchos murieron, pero no porque no tenían los genes adecuados, sino porque no tenían la cultura. Pero apenas se pone a conversar con la gente allá y aprende cómo hacer un iglú, cómo cazar y eso, entonces la cultura es lo que nos permite sobrevivir en todos los medios diferentes.

Pero tampoco es suficiente para explicar algunas cosas. Por ejemplo, uno ve que los neandertales también tenían cultura. Uno busca y ve que tenían pinturas rupestres, uno ve y los neandertales desaparecieron. Y también tenían entierros, tenían también cultura y se piensa que tenían también un lenguaje bastante importante. Y entonces los neandertales desaparecieron.
\ikjpg{28}
Entonces, ¿qué es lo que tienen los humanos que les permite sobrevivir? Así el éxito del que yo hablo, que tenemos aquí una sala completa llena de gente conversando. No se ven chimpancés llenando una sala nunca.
\newpage
\ikjpg{04}
Y entonces la cosa que tenemos nosotros es que tenemos tecnología. Eso sí, los humanos tienen tecnología, los Homo sapiens, bueno, los Neandertales también eran humanos pero ellos no tenían tecnología.

¿Qué es la tecnología? No es solo herramientas, porque hay animales con herramientas, son herramientas acumulativas, es decir, que tenemos herramientas que nos permiten hacer otras herramientas, que nos permiten hacer otras herramientas y eso se va haciendo, eso es una acumulación histórica.

Por ejemplo, primero tenemos el fuego. Esto es una primera herramienta. Después tenemos la arcilla, que la utilizamos para hacer formas. Lo cocinamos en el fuego, obtenemos cerámica. Con eso podemos fundir metales.

Con los metales podemos hacer herramientas. Y así uno se da cuenta que viene todo a construir. Si uno quisiera hoy hacer, no sé, una tostadora de pan a partir de nada, estando ahí en la selva, sería completamente imposible porque hay ese fenómeno de acumulación.

Entonces ahí lo que vamos a tratar de ver es cómo llegamos a esto, cómo llegamos a ese cohete que va en el espacio. ¿Cuál es el camino que nos lleva ahí?

\ikjpg{05}

Y un elemento importante de entender es lo que se llama la coevolución gene-cultura. Es decir, el hecho de que la cultura va a tener una influencia sobre los genes, y claro, los genes también tienen influencia sobre la cultura, pero lo que es menos conocido es la influencia de la cultura sobre los genes.

Entonces, un ejemplo, uno de los primeros ejemplos que uno puede tener en la historia de la humanidad, o mejor dicho, de los proto-humanos, porque se conoce el fuego, los humanos utilizan los proto-humanos, lo que venía antes de los humanos, utilizan el fuego desde hace más de un millón de años. Ya se tienen las evidencias arqueológicas.

Entonces, lo que hace el fuego es que permite, a partir de un momento en que uno tiene el fuego, alguien que tenía un sistema digestivo pequeño, unos dientes pequeños y un cerebro grande, antes del fuego no podía sobrevivir. Pero una vez que se adopta el fuego, que es un elemento cultural, Yo decido de utilizar el fuego de manera regular, hay un aprendizaje de cómo se utiliza el fuego, cómo se puede volver a prender, cómo se cuida.

A partir de que pasa esa adopción de ese elemento cultural, entonces el cuerpo humano puede cambiar y nosotros somos resultado de eso. Es decir, tenemos unos dientes pequeños, un sistema digestivo pequeño y sin fuego no podemos sobrevivir la comida, no podemos alimentarnos suficientemente, por lo menos los niños cuando crecen no pueden alimentarse suficientemente sin fuego. Y entonces eso permitió el crecimiento del cerebro.

Hay que ver también que el fuego lanza o tiene efectos también sobre más cultura, porque cuando hay un fuego la gente se reúne y tiene interacciones sociales. Puede estar de noche, alarga el tiempo, ante un tiempo que era un tiempo durmiendo, se transforma en un tiempo de comunicación social en el que la gente puede socializar. Entonces eso va a traer más cultura, lo que va a empujar a un cerebro más grande. Entonces aquí tenemos un ejemplo de coevolución en la cultura.

Hay otros ejemplos, como por ejemplo el de la puntería. Al adoptar un método de cacería en que uno le lanza proyectiles a las presas, eso va a hacer que alguien que tenga mejor puntería, es decir, unos tendones, unos nervios, una coordinación visual con el gesto, esa persona va a tener más éxito reproductivo para poder brindarle más comida a su familia. Y entonces con eso se va a desarrollar la puntería. No hay ningún animal que tenga puntería, pero que tenga los gestos de puntería. Aquí sí, tenemos un deporte aquí en Venezuela, que ganaron el campeonato hace poco, de puntería.

Pero esos genes de la puntería vienen de un elemento cultural que se adoptó, que es la caza con proyectiles, y entonces después los proyectiles se pusieron más y más sofisticados, lanzas, flechas, en vez de simplemente piedras probablemente al principio.

Otro elemento que también es cultural es el hecho de sudar. No hay ningún animal que suda por ahí, y eso tiene su lógica porque sudar uno pierde su humedad y se muere de sed. Entonces no parece una buena idea. Pero los humanos pueden sudar porque ellos adoptan un elemento cultural que son los recipientes, que puede ser una cáscara de huevo de avestruz, una tapara por aquí en Venezuela, y al adoptar eso, sudar se transforma en una cosa adaptativa, porque podemos correr, el humano es el único que puede correr durante un maratón, Y entonces es una técnica de caza que es seguir el animal hasta que esté agotado y también es para huir un peligro.

Y entonces, claro, el sudar se transforma en algo, una ventaja evolutiva, pero a partir del momento que uno adopta la cultura del recipiente. Y entonces, bueno, ya estamos de acuerdo que la adopción de la cultura va a empujar los genes y en particular, por ejemplo, la del fuego va a empujar a un cerebro más grande.

\ikjpg{06}

Ahora vamos a lo que nos interesa más, que es el lenguaje.

Entonces el lenguaje aparece, hay pájaros con lenguaje, el lenguaje ya es algo que existe fuera de los humanos. Y aquí tenemos el ejemplo de estos simios que se ha identificado muy claramente, que tienen mínimo tres palabras y entonces como tienen una palabra para águila, una palabra para serpiente y otra para leopardo. Entonces uno se agrava el grito que indica serpiente y entonces cuando se pone el grito enseguida se ponen de pie y miran al suelo a ver si hay una serpiente. Entonces significa que entendieron la palabra. Si lo que decimos es águila, miran al cielo y se esconden debajo de unas matas.

Si se le dice leopardo, suben lo más alto del árbol y miran alrededor de ellos. Entonces significa que tienen ese vocabulario.

Y entonces, claro, ahí el lenguaje tiene una ventaja porque les permite escapar de manera colectiva más fácilmente a los predadores.

\ikjpg{07}
Pero en los humanos el lenguaje toma otro camino.

Entonces aquí tenemos a este Robin Dunbar, que él emite la hipótesis que ahora mucha gente es lo que más se ha adoptado, cuál es, por qué el lenguaje se pone, por qué nosotros somos seres de lenguaje, tanto lenguaje que tenemos.

Y entonces él observa que él toma conversaciones en todos los ámbitos, cazadores, recolectores en el Amazonas, el directorio de una empresa, cualquier tipo de cosas, ve las conversaciones y se da cuenta que el 60-70\% de todas las conversaciones son hablar de otros. No solo para chismear, pero aunque puede ser, mira, uno se da cuenta, la gente cuando habla con sus colegas, y él me dijo esto, y yo le dije aquello, etc. Uno tiene ese tipo de conversaciones. Y el lenguaje se desarrolla, es muy bueno para ese tipo de conversaciones. Uno tiene, por ejemplo, unas estructuras gramaticales en que se puede decir, yo pienso que él cree que él no sabe.

Y ahí hay una cantidad de indirecciones, yo creo que los humanos pueden llegar hasta ocho veces. Uno dice, él piensa que yo pienso que él piensa que yo pienso que él piensa, ocho veces. Pero otro tipo de razonamiento en el lenguaje no lo podemos captar. Elementos de lógica, por ejemplo, que parecieran como que es más útil, no sabemos. Pero para lo que es hablar de las relaciones humanas, el lenguaje tiene mucha riqueza.

Entonces eso dice que el lenguaje lo que tiene es un rol social. Es decir, que el lenguaje tiene el mismo rol, por ejemplo, que se llama aquí grooming, que son, por ejemplo, los chimpancés. Ellos hacen alianzas, hacen amistades, limpiándose la piel, quitándose los piojos, y se dan cuenta que los que se quitan los piojos después se comparten comida cuando hay escasez de comida. Entonces, el lenguaje tiene ese rol. La gente va a hablar como para hacer amistades.

Entonces el lenguaje trae un poder social y tiene una ventaja, una ventaja evolutiva. El que utiliza el lenguaje va a tener como aliados, amigos o va a tener más poder social que le da una ventaja reproductiva.

Entonces el lenguaje tiene tanto poder que va a ir aumentando y aumentando y aumentando. Porque esa es la única razón por la cual uno necesita un cerebro tan grande. Porque cuando empieza al principio no hay tanta tecnología, no hay tantas cosas que aprender, uno no necesita un cerebro muy grande para eso. Pero para lo que es el lenguaje sí hay una inflación en la que aparecen más y más elementos, más y más complejidad.

¿Por qué? Porque eso produce una ventaja reproductiva. Así como la cola del pavo real que se pone más y más grande porque eso produce una ventaja reproductiva. Pero producen una ventaja reproductiva hasta cierto punto, porque si uno tiene el cerebro demasiado grande, eso finalmente no es bueno para la supervivencia. Tenemos por ejemplo el ejemplo de las mujeres, antes de la medicina moderna, la causa de muerte, la cuarta parte de las muertes se producían por parto, porque claro, tienen muchos partos, cada parto con un riesgo, y entonces las mujeres morían de esa razón, y los pobres niños con una cabezota, un cuerpo chiquitito, no pueden hacer nada.

Eso es... Es decir, que el tamaño del cerebro se llevó a un límite tal, por ese desarrollo del lenguaje, que es como el de la cola del pavo real, que llega a un límite tal, tiene una ventaja reproductiva, pero llega a un límite que casi es un problema.

Pero también lo que hace es que mientras más se desarrolla el lenguaje, eso permite unos grupos sociales más grandes. Mientras más uno ve que, por ejemplo, los chimpancés tienen unos grupos de 15, 20 personas, a veces hacen una alianza que puede subir hasta 50, y los humanos tienen más bien unos grupos como de 50 personas que pueden tener alianzas de millones y más que eso.

Pero en aquella época probablemente no habían millones.
\ikjpg{08}
Pero ahí es que aparece entonces la primera revolución cognitiva y que es el principio de la humanidad moderna, de lo que podríamos decir el Homo sapiens sapiens. Y eso lo que pasa es que con ese lenguaje que se pone más y más sofisticado aparecen innovaciones culturales que permiten grupos más grandes. Puede ser el hecho de tener clanes, pueden ser historias, pueden ser religiones que unen a la gente, cosas así.

Eso es lo que desarrolla Yuval Harari en su libro Sapiens, que es las historias que uno cuenta, las alianzas que permiten grupos grandes de colaborar, de trabajar juntos. De tener una inteligencia colectiva entonces lo que produce es a esos grupos grandes es que hay una cosa es que va a permitir la acumulación cultural va a permitir la tecnología es decir la construcción de herramientas con sobre herramientas y eso y eso se produce sólo porque uno ve que durante mucho tiempo no pasa nada y de repente hay como un efecto de umbral que pasamos a un cierto nivel y hay una explosión entonces este ahí lo que pasa es que cuando uno llega a cierto cierta cantidad de personas que están porque forman un grupo ahí se puede producir la acumulación cultural imagínate que tengamos un grupo de 15 personas y este y se hacen este que hay una persona que sabe hacer la tal cesta que sirve para la pesca pongamos tal tejido de es esta.

Bueno, hay un accidente o una enfermedad, una cosa, esa persona se muere. Se pierde el conocimiento si no ha tenido tiempo antes de transmitirlo. Pero si hay tres personas que lo hacen porque tenemos un grupo más grande, ahí no se pierde el conocimiento y se puede producir ahí empieza a producirse la acumulación de tecnología, de cultura y entonces ahí empieza la cultura va creciendo y esa Y también esa cultura más grande permite inteligencia colectiva más grande, grupos humanos, que no solo es el grupo inicial de la tribu, pero también este comercio entre las tribus y va creciendo. Entonces eso se ha observado y cuando los grupos humanos se ponen muy pequeños, pierden tecnología, como pasó en Tasmania, que se quedaron cortados de Australia y perdieron casi toda tecnología, quedó lo mínimo de cosas que sabían hacer.

Entonces, a partir de ese momento, hace 50.000 años más o menos, entre 70 y 50.000 años, los humanos se riegan por todo el planeta, llegan a Australia, se cubren. Y también había grupos humanos que estaban aquí y allá, como los Neandertales, posiblemente otros sapiens que estaban ahí. Bueno, estos grupos que vienen con esas innovaciones culturales que les permite tener inteligencia colectiva, que les permite crecer la cultura, desaparecen a todos esos otros humanos en muy poco tiempo. En América no llegan de inmediato hace 50.000 años porque hay que esperar que la época glacial termine. Hay un paquete de hielo al norte de América que impide entrarle a América, pero terminan llegando a todo el planeta.

El último lugar que llegan es a Madagascar. Los humanos, los sapiens.
\ikjpg{09}
Entonces aquí ya podemos atacar las cosas más conocidas, que son las otras revoluciones. Son más conocidas también porque hay historia. Ahí no estamos trabajando con hipótesis, ya tenemos escritura.

Entonces la segunda revolución cognitiva es la de la escritura. Es la de la escritura. Entonces, ¿cómo empieza la escritura? Empieza primero con que hay la agricultura. A partir del momento que hay agricultura va a haber excedentes de producción de alimentos y entonces se produce dominación de una élite sobre los campesinos y entonces tienen impuestos.

Y los impuestos entonces se forman con la necesidad de la contabilidad. Y eso aparece con la contabilidad y entonces la escritura cuneiforme, aquí de este lado en Asiria, lo que es Irak ahora, empezó como hace 5.500 años, pero empezó con la contabilidad.

Y eso permite, de tener la escritura, permite la creación de ciudades. Las ciudades, todas las ciudades existen con escritura, a lo mejor está la excepción aquí de los incas, y esto es el kipú, que en realidad ellos solo tienen la, bueno, no se sabe, no se ha descodificado todo, es posible que también haya un poquito de escritura, pero más que todo es una forma de contabilidad, en que se van con unos nudos, se va contando, entonces sí pudieron hacer la ciudad, y no forzosamente llegaron a lo que vino con la escritura, que fue más allá de eso. Por ejemplo, a partir del momento se empezó con la contabilidad, pero después empezó a escribirse ficción, se empleó para muchas otras cosas la escritura, y eso es lo que llamamos las civilizaciones, nacen de esa escritura, que son unos grupos más grandes que pueden acumular más cultura, y también tienen ya esa información que les permite escribir, entonces ahí nace ciencia, tienes la ciencia en Europa, en el Mediterráneo, en la India, en China, Aparecen ciencias y también hay tecnologías agrícolas.

Por eso es que es una exposición cognitiva que nace de la escritura.

\ikjpg{10}
La siguiente revolución cognitiva es la de la imprenta, Gutenberg, del título.

Y esa aparece en Europa a mediados del siglo XV, más o menos cuando Cristóbal Colón llega a América, y eso produce cambios en Europa, en particular sobre la religión. Allí está la reforma protestante, no que sea tan importante en sí, pero es interesante en la historia de la imprenta, porque los únicos libros que se imprimían, que no se imprimían, que existían en aquella época, era la Biblia. Más o menos ese era el único libro que se escribía y entonces en el catolicismo los únicos que leían eran los curas y ellos leían el libro a los demás. Porque es que los libros habían muy poco, para hacer un libro se necesitaban años de trabajo, los copistas que copiaban esos libros y costaba muy muy caro fabricar un libro. Era un objeto de lujo absolutamente impresionante.

Pero cuando aparece la imprenta lo primero que se imprime es la Biblia y entonces aparece ahí una rama de los cristianos que dice que uno tiene que leer uno mismo tiene que leer la Biblia la gente tiene que tener un libro lo tiene que leer pero también es un movimiento político y entonces se imprime en panfletos en los primeros 50 años de la imprenta 80\% de todo lo que se imprimía tenía que ver con los protestantes. O era Biblia, o eran panfletos, entonces todo empezó con esa cosa de la religión.

Pero después de todo ese conocimiento se empezaron a imprimir libros de otro tipo, antiguos, que venían del mundo árabe, y ahí entonces eso lanzó la revolución científica y también la gente empezó a imprimir todo tipo de libros, eso se transformó en un negocio.

En realidad, las primeras fábricas, aquí tenemos costumbre de tener objetos que son todos idénticos, que salen de una fábrica. Bueno, los primeros objetos que hemos fabricado así son los libros. Esas eran las primeras fábricas y entonces ataban. Nació ahí el capitalismo industrial. Nace mucho de la imprenta y se va regando.

Entonces, de la revolución industrial nace, de la revolución, es decir, científica, nace la revolución industrial y entonces esto produce todo el, todavía estamos en eso, ese es el boom que tenemos de toda la riqueza con un aumento de población. También lo vimos que hubo cuando la primera revolución cognitiva, cuando sale Out of Africa, hay una multiplicación de la cantidad de gente por mucho. Lo de las ciudades también produjo una multiplicación de la cantidad de gente con la agricultura y las ciudades, y esto de nuevo produce esta explosión que tenemos de miles de millones de personas.

\ikjpg{11}
La imprenta también nos trae un impacto político muy grande. Una de las cosas que nace en ese momento es el Estado-Nación, los idiomas por países.

Antes lo que había es que toda la gente culta, la gente que leía, leía en latín o en griego, y la élite intelectual, la poca gente que leía, leía en esos idiomas.

Pero cuando aparece la imprenta, entonces aparecen idiomas como el francés, como idiomas que van tomando el inglés, que van tomando un espacio. Y entonces se forma esa noción de que hay una zona territorial con un idioma, una nación, nace el nacionalismo en ese momento. Antes las fronteras se movían, había guerras varias por siglo, y 10 años la frontera estaba por aquí, después estaba por allá, y los duques y los reyes y todo venía subdividido. En ese mundo medieval, ahora tenemos el Estado-Nación, eso se fue construyendo poco a poco. Italia, Alemania, esos países se unificaron en el siglo XIX.

Aquí en América también los países nacieron en el siglo XIX. Y en África, hasta en el siglo XX, aparecieron países que todavía tienen su... Es decir, el nacionalismo es un invento, es algo que nace con la imprenta. Está el libro de Benedict Anderson, que habla muy bien del nacimiento del nacionalismo. Pero eso también trae la democracia, porque tienes educación universal, todo el mundo aprende a leer, tienes la prensa, la prensa escrita, que es como un contrapoder contra los gobernantes, el imperio de la ley.

Antes la ley no se podía aplicar porque simplemente los textos de leyes no se podían difundir. Eso ahora sí se puede. Entonces antes la ley la hacía una persona que decretaba esto y había muy pocos textos de leyes porque simplemente no se podían producir.

Entonces eso cambia y eso es también la imprenta. Ese es un cambio completo del mundo, con un cambio de poderes, caen los reyes, viene la democracia. Esto es bastante reciente, todo esto. Todavía estamos en ese movimiento que impulsa la imprenta.

En China la imprenta llegó antes, llegó como en el siglo VII, y en China también produce un efecto muy importante de la unificación del imperio. Antes había reinos, no había unificación. Eso ayuda a la unificación del imperio, tiene una unidad cultural, tiene también el invento de la meritocracia. Los mandarines, los que mandan ahí, tienen un concurso en que tienen que saber leer, escribir y eso lo produce la imprenta.

Y por cierto, mucha de la revolución científica industrial viene de China, que ya ellos habían tenido un principio de eso, la pólvora, la propia imprenta, claro, viene de China. Muchas innovaciones vienen de China, pero la diferencia es que en China la imprenta está en manos del poder, el emperador es el que imprime, no todo el mundo. Entonces no nace por ese poder centralizado, no nace como pasó en Europa, en que había muchos poderes, contrapoderes, la Sorbonne prohíbe la imprenta en Francia. Los pensadores se van, se mudan a los países vecinos y entonces todos tenían esa dinámica que en China no existió, no se produjo esa dinámica como pasó con la imprenta en Europa.
\ikjpg{12}
Entonces seguimos aquí en nuestro camino en la información y en el siglo XIX aparece el telégrafo.

Claro, muchas cosas, pero ahí quiero resaltar el telégrafo porque vemos que ya al final del siglo XIX tenemos una red de telégrafos por todo el planeta. Uno puede mandar un mensaje que llega bastante rápido, bueno, sí toma algunos días porque una persona tiene que reenviarlo y eso, pero uno puede mandar un mensaje bastante rápido de América a Australia, de cualquier parte. Esta es la mundialización que conocemos ahora. Claro, ahora no son esos mismos cables que antes eran barcos de vela que iban lanzando los cables, pero ahora son fibras ópticas que hacen eso.

Pero en realidad, el Internet empieza en ese momento.

Y por cierto, en aquel momento, no sé si ven aquí a la izquierda la imagen, este es un primer teclado para mandar mensaje a distancia. Uno escribe aquí directamente las letras, en vez de utilizar el morse, uno puede escribir directamente las letras.

Entonces, claro, tomaron como base un piano, pero uno escribía, se escribía y eso es el ancestro, eso se llamaba un teletipo, teletype, así como el teletype que se utilizó en las computadoras, que es lo que vemos al terminar Linux, que viene cuando uno escribe que es un teletipo, bueno, ese es el teletipo, el ancestro del teletipo es esto y es en el siglo XIX y es como un piano.

Entonces, está aquí el telégrafo, pero también entonces aparece la computación.
\ikjpg{13}
Aparece la computación, pero la computación en aquella época, lo que se llama una computadora, es una persona. Mayormente mujeres, aunque uno se mira en la imagen ve que también hay hombres, pero mayormente son mujeres. Una computadora es una mujer. Usted ve aquí escrito, dice Computing Division. Esto es una sala de computación, principio del siglo XX.

Una sala de computación, habían así todas las grandes empresas, tenían eso para procesar los datos que tenían, los gobiernos, los bancos, tenían así computación. Con lo que se llamaba una computadora, era una persona en ese momento.
\ikjpg{14}
En los años 40 aparece la primera computadora electrónica que se ve ahí, ese es el ENIAC. Claro, esa computadora no funcionaba muy bien, funcionaba como dos horas, era muy lento y tal. Y en realidad, las computadoras humanas solo desaparecieron en los años 60. Hasta el final de los 60 todavía existían en las empresas, en las universidades, existían computadoras.

Cuando uno decía computadora, por lo menos en inglés, no sé exactamente en Venezuela cómo se decía, pero en inglés, computer era una persona cuando uno decía computer. Pero hasta, fue como en los años 60 que ya uno no tenía que explicar la cosa.
\ikjpg{15}
Entonces a partir de ahí las cosas van muy rápido, pero uno se da cuenta que hay una aceleración de la historia, y en el 69 tenemos internet, en 1969 tenemos internet, ahí son las cuatro primeras computadoras de internet, tres en California, otro en un estado vecino, y esto es una computadora IBM 360 de la época, una de las cuatro computadoras, esta aquí es la 360, y esto no es la computadora, esto es el teclado únicamente, esto es el teletipo, el mismo teletipo del piano con las letras, no hay pantalla todavía, el texto se imprime a medida que lo produce la computadora, y la computadora per se ocupa toda una sala. 
\ikjpg{16}
Entonces las cosas van ahí muy rápido porque 10 años después tenemos las computadoras personales, las computadoras que uno puede tener en su casa, que tienes el Apple II, que ese es del año 1980. Yo tenía uno así aquí en Caracas y exactamente ese modelo.

Y ahí es que la gente empieza a tener computadoras en su casa, lo que parecía imposible tener computadores en su casa. Y además las computadoras no existían. 30 años antes no había computadoras. Y esto es una computadora de bolsillo con conexión a internet. Esto tiene 20 años.
\ikjpg{17}
Tiene 20 años. Es el iPhone. Bueno, no puse todo de la computación y el web, cantidad de innovaciones, pero ustedes se dan cuenta que las cosas van muy rápido.

Y vamos a ver la inteligencia artificial la inteligencia artificial se emplea hay personas que hablan de eso ya en los años 40 como John Von Neumann que en Princeton que es de los que inventó las computadoras modernas empieza a hablar de inteligencia artificial en ese momento también Alan Turing empieza a hablar pero en los años 50 hay una conferencia donde un tal McCarthy habla de inteligencia artificial y se empieza la investigación en inteligencia artificial.

Aquí vemos lo que decía este investigador en el 58 en inteligencia artificial. Él dice, dentro de 10 años una computadora digital será campeón de ajedrez en el 58.

Aquí hay dos cosas interesantes. Él no dice una computadora, dice una computadora digital.

¿Por qué? Por lo que yo decía, es que una computadora es una persona. Computadora digital.

Por ejemplo, Shannon, el de la teoría de la señal, él se casó con su computadora. Era muy tímido y la única persona que le hablaba todos los días era con su computadora. Y después de 10, 20 años juntos, dijeron, bueno, vamos a casarnos. Y eso fue, él se casó con su computadora. Bueno, entonces ahí en aquella época uno decía, y estaba esa angustia de que las computadoras digitales iban a reemplazar a las computadoras humanas, que es como algo parecido a las angustias que tenemos ahora con la inteligencia artificial.

Claro, entonces ahí él se equivocó y dijo dentro de 10 años. No fueron 10 años, fueron 40 años. Pero tenía razón. El campeón de ajedrez terminó siendo una computadora digital. Se tardó 40 años.

Después tenemos aquí en el 65, en el 70, otros investigadores conocidos de la historia de la inteligencia artificial que dicen que dentro de 20 años una máquina será capaz de hacer cualquier trabajo que hace un humano. Eso todavía no nos lo están prometiendo ahorita. No ha pasado. Y Marvin Minsky decía que de 3 a 8 años tendremos una máquina con la inteligencia general de un humano.

Una cosa también ahí que es muy interesante cuando uno habla de inteligencia artificial, es que siempre, siempre nos está comparando con los humanos. El tema de la inteligencia artificial es que en realidad no tiene tanto que ver con la máquina, tiene que ver con los humanos.
\ikjpg{18}

Pero hay algo un poco extraño, porque si cuando nosotros hablamos de un avión, claro que los aviones fueron inspirados primero por las aves, pero nosotros no comparamos un avión con un alcatraz. No estamos diciendo que un avión vuela tan bien como un alcatraz, pero cuando hablamos de inteligencia artificial, sí lo comparamos con los humanos como si los humanos fueran lo más inteligentes que hay. Pero no es verdad, porque hay diversos tipos de inteligencia. Por ejemplo, las ardillas, durante el verano esconden nueces, hasta 10.000 nueces las esconden en un huequito, y se acuerdan después cuando llega el invierno, saben dónde están y las buscan. Ningún humano puede acordarse de donde ponen 10.000 cosas.

Uno pone tres cosas y ya no las consigue en su casa. Porque no tenemos ese tipo de inteligencia, como lo vimos, el tipo de inteligencia que tenemos es una inteligencia social, no es una inteligencia general, es una inteligencia social que nos permite colaborar en grandes grupos. Pero aquí tenemos esa tendencia de comparar con los humanos.
\ikjpg{19}
En realidad es hasta peor que eso. Uno no debería decir inteligencia artificial, uno debería decir inteligencia humana artificial. Y lo que llamamos inteligencia artificial es siempre cuando la computadora está al punto de hacer algo que nosotros los humanos hacemos, que nos entra una angustia.

Entonces nosotros pensábamos que a jugar a ajedrez, eso era una actitud humana.

Cuando en los años 90 cuando Deep Blue le gana a Kasparov, los titulares de los periódicos dicen, si tu trabajo se parece a jugar a ajedrez puedes estar inquieto claro entonces la gente decía, si de verdad lo que yo hago es más es hasta más fácil que jugar a ajedrez voy a estar muy inquieto, entonces claro uno definía el jugar a ajedrez es como una actitud humana muy fuerte entonces viene Deep Blue le gana completamente a Kasparov y después todavía más potente y eso ¿quién dice que ahora jugar a ajedrez es una algo, es un signo de inteligencia humana? Nadie, porque las computadoras son mucho más fuertes que eso, entonces eso es inteligencia mecánica, eso no nos interesa, eso no es inteligencia humana artificial, esto es inteligencia mecánica, siempre lo fue se nos olvida que hace poco consideramos que eran los humanos que hacían eso, se nos ha olvidado y después hace tanto tiempo que calcular una raíz cuadrada era una cosa que hacían los humanos.

Ya nos parece, ah, calcular más raíz cuadrada, eso lo hacen las máquinas, eso no nos concierne, nosotros no estamos preocupados por eso.

Entonces, claro, cuando estamos ahorita que de repente las computadoras están al punto de ser tan buenos como nosotros en la conversadera, en el chatbot, de repente nos inquietamos. De repente nos inquietamos y decimos, estoy ahí, ¿qué va a pasar? ¿Esto es una tragedia? ¿Qué es lo que está pasando?

Pero en realidad, posiblemente, no sé si es la verdad, pero si de repente los robos se ponen y son igual de buenos que nosotros en conversar, en este chiste, ya no vamos a considerar eso como inteligencia artificial. Y ya está pasando. Hay gente que está diciendo, no, ahora imagínate, un niño de 8 años, uno le dice, recoge la mesa y friega los platos. Va y lo hace. Eso en una computadora no lo puede hacer.

Entonces, ah, no, eso es la inteligencia humana. Esa es la inteligencia humana artificial. Entonces ahí viene ese proceso que es a la frontera cuando las máquinas empiezan a llegar a una capacidad que consideramos como humana. Y entonces ahí mismo la descartamos y decimos no, no, eso no es humano, eso es mecánico. Claro, porque si lo hace la computadora es que es mecánico.

Eso es la inteligencia artificial, es eso.
Hay que acordarse por ejemplo que reconocer una cara, eso era considerado hace 10 años, en los centros de investigación en inteligencia artificial, uno trataba de hacer como hacer para reconocer una cara, eso parecía una cosa, bueno, solo los humanos pueden hacer eso. Ahora uno desbloquea su teléfono mirando, lo reconoce la cara de uno, y además las computadoras son mucho más fuertes que uno, un humano puede reconocer 100 caras, 200 caras, pero una computadora va a reconocer millones de caras, porque en un aeropuerto van hasta las cámaras, identifica a todo el mundo. Ya las computadoras sobrepasaron completamente a los humanos en eso, lo que hace 10 años no lo era y se llamaba inteligencia artificial y ahora nadie se preocupa, no hay inteligencia artificial. A lo mejor había preocupación hace un tiempo con eso, pero ahora no pareciera.

Entonces, bueno, ahí está a lo mejor es un poco esquemático lo que estoy diciendo, pero ahí está ese tema de que la inteligencia artificial tiene también que ver con nosotros, cómo nos definimos relativamente, como decía yo al principio, a los animales y ahora relativamente a las computadoras.

Pero en realidad, fuera de comparar con los humanos, de la rama específica de la computación en que se trata de imitar a los humanos. 

\ikjpg{20}

Hay una cosa que sí, que tiene más fuerza, diría yo, que eso, y que explica por qué siempre nos sorprende todo.

Lo que tenemos, ustedes han oído hablar de la ley de Moore.

Gordon Moore, el fundador de Intel, notó en los años 60 que todos los años o todos los dos años, uno multiplicaba por dos la cantidad de transistores que se ponían en una área dada. Y eso hacía que, claro, como los transistores estaban más cerca y las cosas iban más rápido, y también había más potencia de cómputo en la misma superficie. Entonces había una multiplicación de la potencia de cálculo que había en cada chip.

Bueno, eso más pequeño, más pequeño los transistores, está el límite del tamaño de los átomos, uno no puede ser más pequeño siempre, así de manera exponencial, el doble todos los años. Y pasó que como en el 2010 o en el 2015, eso dejó de ser verdad. No se podían poner más transistores en una cosa pequeña.

Pero la ley que sí, si uno lo pone de otra manera, Kurzweil lo pone de esa manera. Él dice cuánto capacidad de cómputo puedo comprar con mil dólares. ¿Cuál es la capacidad de cómputo que puedo tener? Y esa sigue una forma exponencial, esta es una escala logarítmica, si una línea recta sería exponencial, pero ahí uno ve que va hasta subiendo, es decir que es más que exponencial. Y va más rápido que una exponencial, pero bueno, es más o menos exponencial.

Y entonces, ahí eso sigue así, eso sigue más allá del 2015, sigue andando, uno se da cuenta que sigue habiendo más y más computadoras. Y entonces lo que él dice es que en el 2029, aquí en el 2029, tendremos por mil dólares, el 2029 es dentro de poco, de por mil dólares tendremos la misma capacidad de cómputo que un humano. Mira, pareciera que nos estamos muy lejos, es factible.

Claro, la capacidad de cómputo no es todo, uno también necesita como que el software, pero por el momento todas esas predicciones han sido más o menos, se han ido acertando. Claro, entonces la potencia, eso es lo que explica la sorpresa que tenemos, Porque exponencial es como muy rápido. Si es una multiplicación por dos cada año, eso significa que en 10 años es una multiplicación por mil. Y en 20 años es una multiplicación por un millón. Es decir, es enorme.

Tenemos unas computadoras que son un millón de veces más potentes que hace 20 años. Entonces siempre estamos como con la sorpresa. Nos viene algo y nos dice que las computadoras pueden hacer eso. Pero es porque viene eso. Entonces él dice que en el 2049 las computadoras por mil dólares uno tendrá la misma potencia que todos los humanos reunidos.

Entonces bueno, eso es ahí como una religión que se forma de la singularidad, que en el 49 llega el apocalipsis, no sé cuánto. Pero aparte de eso, es verdad eso del crecimiento exponencial de la potencia de cómputo.
\ikjpg{21}
Entonces, como ya no está la ley de Moore, que se pone más y más potencia en cada chip, ¿cómo sigue creciendo la capacidad de las computadoras?

Bueno, lo que llamamos los data centers.

Entonces, aquí tenemos los componentes adentro de un teléfono o dentro de una computadora. Tenemos los chips, así. Y esto es más o menos la misma cosa. Pero aquí, esto son unas carreteras son árboles, son esto pero ustedes ven, estos son parecidos a los chips pero las computadoras se van poniendo más y más grandes y están ocupando unos espacios más y más grandes pero la cantidad de computadores sigue aumentando sigue aumentando, y entonces claro, también se pone más grande, también necesitan más energía para para funcionar, entonces hay un aumento de la cantidad de energía que necesitamos para ese cómputo se estima, por ejemplo, en California ahora ya estamos al 20\% de la energía eléctrica que se utiliza para los data centers y a nivel de planeta va a ir creciendo hasta un momento que yo diría el 80\% de toda la energía que vamos a utilizar en el planeta será para esto.
\ikjpg{22}
Entonces no supe exactamente cómo representar la energía, pero aquí puse para que uno se dé cuenta que si hay mucha energía por detrás, estos son los tubos de enfriamiento de un data center de Google hace unos años pero se sigue necesitando enfriando y ustedes se dan cuenta el tamaño aquí de estas tuberías la cantidad de eso es para enfriar las computadoras es inmenso es una cantidad de energía inmensa, entonces lo que pasa es que todo ese aumento de energía produce un límite es decir que ahora ya hay unos data centers que dice no, no queremos data center ahí porque eso está calentando el río y se mueren los peces eso es una polución térmica local pero si sigue creciendo así ya va a ser al nivel del planeta el calentamiento que va a producir toda esa cantidad de computadoras que va creciendo acuérdense de manera exponencial una multiplicación por mil en 10 años, es decir que dentro de 20 años vamos a tener un millón más de data center de lo que tenemos ahora algo así, más o menos entonces claro, uno se da cuenta que eso va a seguir así parece imposible bueno, 

\ikjpg{23}
ahí es que viene la etapa siguiente en el 2025 se empezó a hablar de las computadoras espaciales las computadoras espaciales entonces, ¿por qué?

Porque también mandar cosas al espacio se ha puesto más barato aquí tenemos el cohete Starship de SpaceX la compañía de Elon Musk ¿por qué se abarata? Porque hay una capacidad de cómputo muy importante que permite a los cohetes volver a aterrizar bueno, eso no es inteligencia artificial porque un humano no lo puede hacer entonces ya sabemos que eso no importa pero ahí tenemos la capacidad de cómputo que está haciendo algo que nos permite que no lo hacen los humanos y entonces tenemos Google Suncatcher ellos ponen un artículo que ellos tienen y en que van diciendo que van haciendo un plan para hacer un data center en el espacio. Es decir, una cantidad de satélites conectados que hablan con láser. Si tú googleas Suncatcher, ahí lo vas a conseguir.

Y ellos van a lanzar su primer satélite en el 27.

Tienes una empresa, una startup pequeña que ya pidió permiso para hacer 88.000 satélites en el espacio para montar su data center en el espacio y ya lanzaron su primer cohete en noviembre del año pasado.

Blue Origin. Eso es de este año que ellos anunciaron esto.

Blue Origin es la compañía de Jeff Bezos, el fundador de Amazon, y ellos pidieron un permiso para lanzar 51.000 satélites. El primer satélite lo lanzan este año.

Space X, la compañía de esas que hace los cohetes de Elon Musk, compró la compañía de Elon Musk, que es el equivalente de OpenAI, que se llama xAI, y ahí se fusionaron las dos compañías, y ahora el negocio de eso es hacer data center en el espacio. Y ellos pidieron permiso para un millón de satélites, y se piensa que los primeros se van a lanzar en el 28.

\ikjpg{24}

Entonces todo eso parece ciencia ficción, pero no lo es tanto. Esto es en el 20, no conseguí imágenes más recientes, pero en el 2023 esto es Starlink. Ustedes se acuerdan que hace, pongamos, 5 años Starlink no existía. Ahora hay 10.000 satélites, 10.000 satélites. Y uno va aquí a Multimax y compras una antena Starlink y ya tienes tu conexión que cuesta 30 dólares mensuales.

Es decir, que esto ya existe.

Esto está hecho para la comunicación, pero imagínate que en cada satélite tiene ahorita una pequeña capacidad de cómputo, pero le pones más capacidad de cómputo y ya lo tenemos. Entonces eso parece ciencia ficción, un millón de satélites de aquí, pero no lo es tanto.

Entonces, si tú me dices cuál es el futuro de la computación, yo te digo es esto. Es decir, aquí tenemos, a lo mejor no va a ser 10 años, porque las razones económicas van a ser 20 años, un poquito como este que dice que el campeonato de ajedrez lo iba a ganar una computadora digital, bueno, a lo mejor no va a ser ahorita de inmediato, pero de qué va a pasar, yo creo que sí.

Y entonces esto, y ahí no hay límite para el crecimiento exponencial, porque hay energía solar, una cantidad más o menos ilimitada, y en el sistema solar, bueno, falta mucho, mucho crecimiento exponencial para que no haya suficiente. Y imagínate una fábrica en el espacio utilizando los materiales de la luna para hacer esta cosa, entonces uno va a ver una multiplicación así de las computadoras. Entonces claro, ahí uno está pensando en la inteligencia artificial, pero esto es una angustia muy local, muy de los humanos. Si uno lo mira a una escala más grande, uno se da cuenta que las cosas tienen mucho más empuje que el tema de los chatbots.
\ikjpg{25}
Bueno, entonces, en conclusión, esta es la cuarta revolución cognitiva y estamos en medio de eso. Estamos en medio de eso y yo le puse de la información mágica porque estaba tratando de ver qué es lo que une todo eso y es un poco difícil, porque como estamos todavía en medio de la revolución de la cosa que está pasando, es como difícil tener un análisis histórico de la cosa.

Lo vimos en el siglo XIX, aparece el telégrafo, la fotografía, gramófono, hay todas unas cosas del cine, el teléfono, calculador electromecánico, hay toda una revolución de la información que empieza en el siglo XIX, en el siglo XX tenemos la computadora electrónica, internet, la inteligencia artificial, las computadoras de bolsillo, los satélites de telecomunicación, y bueno, esto sigue andando, esto no se para ahorita mismo.

Entonces, ¿cuál es el futuro? Hay gente que está muy inquieta, que dice los humanos van a desaparecer y los van a reemplazar las máquinas.

Eso es lo que yo puse aquí, obsolescencia del wetware. El wetware es un concepto un poco, vamos a decir, humorístico, en que uno tiene el software, el hardware y el wetware. Es decir, el wet de húmedo, el material húmedo. Nosotros somos el material húmedo de la inteligencia.

Entonces, yo no pienso, yo pienso que sí va a conseguir un lugar, aunque claro, lo otro como va creciendo exponencialmente, la parte húmeda va a ser como más pequeña comparado a todo lo otro, sobre todo si empieza a llenar todo el sistema solar. Si posiblemente vamos a estar ahí y es lo que se llama esa es la otra hipótesis que es el cyborg o el cyborg colectivo, porque nosotros formamos un colectivo, pero el cyborg es esa idea de que nos vamos integrando con las máquinas y eso es más o menos lo que pasa antes uno para acceder a una computadora uno se montaba en un carro y te hacía tantos kilómetros y te acercabas a la computadora Después la computadora llegó a tu casa, después la computadora está en tu bolsillo y ahora están los lentes esos conectados, los tienes así y ya uno no puede tener hasta implante. Es decir, está ese tema de la integración siempre más apretada a la computadora que se va acelerando.

Va acelerando y eso posiblemente eso estamos que nos estamos dirigiendo a eso eso lo que estamos viendo lo que se puede observar claro siempre habrá una gente que va a resistir que no quiere y eso pero habrá un grupo importante de gente que sí si está entrando en ese en esa cosa bueno ahí no está de como difícil hacer las predicciones porque porque la cosa es que viene con una Las cosas se van acelerando.

La primera revolución cognitiva fue hace 55.000 años, más o menos. Después, 50.000 años después, aparece la aparición de la escritura. 500 años después está la imprenta, aparece la imprenta y después, ¿cuánto es? 5.000 años la imprenta, 500 años estamos en esta revolución. Entonces, ¿cuál es el próximo ciclo?

A lo mejor ya estamos en la quinta revolución cognitiva, no nos hemos dado cuenta. Pues serían 50 años el siguiente, si uno toma esa aceleración exponencial.

Bueno, el hecho es que ese es el tema más bien que uno tiene que ver, es el tema de la aceleración.

Más allá de esta cosa de los LLM que quedamos sorprendidos, esos solo tienen tres años y medio y ya todo el planeta los ha utilizado, ya parece una cosa normal. Yo lo que pienso es que vamos a estar de nuevo sorprendidos puede ser este año o dentro de poco por ese tema de la aceleración y yo no creo que la aceleración ahorita que eso se vaya a parar claro, se puede pensar que en algún momento se va a frenar pero no de inmediato entonces estamos en ese tema de la aceleración entonces esa es más o menos la conclusión.

\end{document}